%% toc: P2 22 22.8
\section{Proofs for the comparison of \(L_{\omega_1\omega}\) with \(L_{\infty\omega}\)}
\label{sec:annex-p02-w02-p02-039-lomega-treatments}

Throughout, the notation is that of
Convention~\ref{def:p02-w02-p02-039-notation}.  Finite partial isomorphisms,
potential isomorphisms and their forth and back conditions are those of
Definition~\ref{def:p08-bvb501-potential-isomorphism}: a finite partial
isomorphism is an injective finite map that preserves and reflects every
atomic formula, and the forth condition for a member \(f\) of a potential
isomorphism \(I\) and a point \(a\) supplies some \(g\in I\) with
\(f\subseteq g\) and \(a\in\operatorname{dom}g\); the back condition is the
same with the range in place of the domain.  The metatheory is ZFC.

\begin{annexitem}{annex:p02-w02-p02-039-formula-set}{thm:p02-w02-p02-039-formula-set}{The formulas of \(L_{\omega_1\omega}\) over a fixed signature form a set}
\begin{proof}
Let \(\mathrm{Var}\) be the countable set of variables and let
\(\mathrm{At}\) be the set of atomic formulas over \(\tau\).  It is a set:
terms are finite sequences built from the set \(\tau\cup\mathrm{Var}\) and
finitely many auxiliary symbols, finite sequences from a set form a set, and
an atomic formula is a relation symbol of \(\tau\) or equality followed by
finitely many terms.  Write \(\FV(\varphi)\) for the free-variable set of
\(\varphi\).

Define sets \(F_\xi\) for ordinals \(\xi\leq\omega_1\) by recursion.  Put
\(F_0=\mathrm{At}\).  Let \(F_{\xi+1}\) be the union of \(F_\xi\), of the
results of applying each first-order formation rule (negation, each binary
connective, and each quantifier with each variable of \(\mathrm{Var}\)) to
members of \(F_\xi\), and of all \(\bigwedge\Phi\) and \(\bigvee\Phi\) for
countable \(\Phi\subseteq F_\xi\) with \(\bigcup_{\psi\in\Phi}\FV(\psi)\)
finite.  At a limit \(\lambda\) put \(F_\lambda=\bigcup_{\xi<\lambda}F_\xi\).

Each \(F_\xi\) is a set, by induction on \(\xi\).  At successors, the new
members are images under fixed set operations of members of the sets
\(F_\xi\), \(F_\xi\times F_\xi\), \(\mathrm{Var}\times F_\xi\) and of the
subset of the power set \(\mathcal P(F_\xi)\) cut out by the countability
and free-variable conditions; Power Set, Separation and Replacement make
the result a set.  At limits, Replacement and Union apply.  By construction
\(F_\xi\subseteq F_\eta\) whenever \(\xi\leq\eta\).

\emph{Every member of \(F_{\omega_1}\) is a formula.}  By induction on
\(\xi\): atoms are formulas, and each member added at a successor stage is
obtained by a formation rule of Definition~\ref{def:p02-lomega-one-omega}
from formulas, with that rule's side condition checked at the time of
addition.

\emph{\(F_{\omega_1}\) is closed under every formation rule.}  Since
\(\omega_1\) is a limit ordinal, each \(\varphi\in F_{\omega_1}\) lies in
some \(F_\xi\) with \(\xi<\omega_1\); let \(\rho(\varphi)\) be the least
such \(\xi\).  If \(\varphi,\psi\in F_{\omega_1}\), both lie in
\(F_\eta\) for \(\eta=\max(\rho(\varphi),\rho(\psi))<\omega_1\), so each
first-order compound of them lies in \(F_{\eta+1}\subseteq F_{\omega_1}\).
Now let \(\Phi\subseteq F_{\omega_1}\) be countable with a finite total
free-variable set.  The set \(R=\{\rho(\varphi):\varphi\in\Phi\}\) is a
countable set of countable ordinals.  In ZFC the ordinal \(\omega_1\) is
regular: a countable union of countable sets is countable, which uses the
axiom of choice to select an enumeration of each set.  Hence
\(\eta=\sup R\) is countable, so \(\eta<\omega_1\).  Then
\(\Phi\subseteq F_\eta\), and the conditions on \(\Phi\) put
\(\bigwedge\Phi\) and \(\bigvee\Phi\) into \(F_{\eta+1}\subseteq
F_{\omega_1}\).

\emph{Every formula belongs to \(F_{\omega_1}\).}  A formula is obtained
from atoms by a well-founded tree of applications of the formation rules.
By induction along that tree, using that \(F_{\omega_1}\) contains
\(\mathrm{At}\) and is closed under every rule, the formula lies in
\(F_{\omega_1}\).

Therefore the class of \(L_{\omega_1\omega}\)-formulas over \(\tau\) equals
the set \(F_{\omega_1}\).  Its sentences form the subset
\(\{\varphi\in F_{\omega_1}:\FV(\varphi)=\varnothing\}\), a set by
Separation.
\end{proof}
\end{annexitem}

\begin{annexitem}{annex:p02-w02-p02-039-rigid}{thm:p02-w02-p02-039-ordinal-rigidity}{Ordinals are rigid under potential isomorphism}
\begin{proof}
Each \(f\in I\) is injective and preserves and reflects the atomic formula
\(x<y\), so for \(a,a'\in\operatorname{dom}f\) we have \(a<a'\) if and only
if \(f(a)<f(a')\).

\emph{Claim.}  For every ordinal \(\gamma\), every \(f\in I\) and every
\(a\in\operatorname{dom}f\) with \(\min(a,f(a))=\gamma\), we have
\(f(a)=a\).  The proof is by transfinite induction on \(\gamma\): assume the
claim for all ordinals below \(\gamma\), and suppose toward a contradiction
that \(f\in I\), \(a\in\operatorname{dom}f\), \(b=f(a)\),
\(\min(a,b)=\gamma\) and \(a\neq b\).

\emph{Case \(a<b\).}  Then \(\gamma=a\), and \(a<b<\beta\) gives
\(a\in\beta\).  By the back condition there is \(g\in I\) with
\(f\subseteq g\) and \(a\in\operatorname{ran}g\); choose \(c\in\alpha\)
with \(g(c)=a\).  Since \(f\subseteq g\), \(g(a)=b\).  From
\(g(c)=a<b=g(a)\) and reflection of \(<\) we get \(c<a\).  Then
\(\min(c,g(c))=\min(c,a)=c<\gamma\), so the induction hypothesis applied to
\(g\) and \(c\) gives \(g(c)=c\), that is \(a=c\).  This contradicts
\(c<a\).

\emph{Case \(b<a\).}  Then \(\gamma=b\), and \(b<a<\alpha\) gives
\(b\in\alpha\).  By the forth condition there is \(g\in I\) with
\(f\subseteq g\) and \(b\in\operatorname{dom}g\); let \(d=g(b)\).  From
\(b<a\) and preservation of \(<\) we get \(d=g(b)<g(a)=f(a)=b\).  Then
\(\min(b,g(b))=\min(b,d)=d<\gamma\), so the induction hypothesis applied to
\(g\) and \(b\) gives \(g(b)=b\), that is \(d=b\).  This contradicts
\(d<b\).

Both cases are impossible, which proves the claim.  Given any \(f\in I\)
and \(a\in\operatorname{dom}f\), apply the claim with
\(\gamma=\min(a,f(a))\); thus every member of \(I\) is the identity on its
domain.

\emph{Equality of the ordinals.}  The family \(I\) is nonempty; fix
\(f\in I\).  For \(a\in\alpha\), the forth condition gives \(g\in I\) with
\(a\in\operatorname{dom}g\); since \(g\) is the identity on its domain,
\(a=g(a)\in\beta\).  So \(\alpha\subseteq\beta\).  For \(b\in\beta\), the
back condition gives \(g\in I\) and \(c\in\alpha\) with \(g(c)=b\); since
\(g(c)=c\), \(b=c\in\alpha\).  So \(\beta\subseteq\alpha\).  Hence
\(\alpha=\beta\).
\end{proof}
\end{annexitem}

\begin{annexitem}{annex:p02-w02-p02-039-proper-fragment}{thm:p02-w02-p02-039-proper-fragment}{\(L_{\omega_1\omega}\) is strictly weaker than \(L_{\infty\omega}\)}
\begin{proof}
Two facts are used repeatedly.  First, every structure \(M\) is potentially
isomorphic to itself: the restrictions of the identity map to finite subsets
of the domain form a nonempty family of finite partial isomorphisms, and the
identity restricted to \(F\cup\{a\}\) extends the identity restricted to
\(F\) with \(a\) in its domain and range.  Hence \(M\models\delta_M\) by
Theorem~\ref{thm:p08-bvb501-potential-isomorphism}.  Second, for ordinals
\(\alpha,\beta\geq1\), the structures \((\alpha,<)\) and \((\beta,<)\) are
potentially isomorphic only if \(\alpha=\beta\), by
Lemma~\ref{thm:p02-w02-p02-039-ordinal-rigidity}.

\emph{Clause 1.}  Let \(\alpha\geq1\) be countable.  Then \((\alpha,<)\) is
a countable structure in the countable signature \(\{<\}\), so by the last
clause of Theorem~\ref{thm:p08-bvb501-potential-isomorphism} there is an
\(L_{\omega_1\omega}\)-sentence whose models are exactly the structures
potentially isomorphic to \((\alpha,<)\).  It has the same models as
\(\delta_{(\alpha,<)}\), so the two are equivalent.

\emph{Clause 2.}  By Lemma~\ref{thm:p02-w02-p02-039-formula-set} applied to
\(\tau=\{<\}\), the \(L_{\omega_1\omega}\)-sentences over \(\{<\}\) form a
set \(S\).  Let \(A\) be the class of ordinals \(\alpha\geq1\) for which some
\(\sigma\in S\) is equivalent to \(\delta_{(\alpha,<)}\), that is, for every
structure \(N\) of the signature, \(N\models\sigma\) if and only if \(N\)
is potentially isomorphic to \((\alpha,<)\).  This condition is expressed by
a formula of set theory, so \(A\) is a class.  For \(\alpha\in A\) put
\[
 T(\alpha)=\{\sigma\in S:(\alpha,<)\models\sigma\}\subseteq S.
\]

\emph{\(T\) is injective on \(A\).}  Let \(\alpha\neq\beta\) lie in \(A\)
and let \(\sigma\in S\) be equivalent to \(\delta_{(\alpha,<)}\).  Since
\((\alpha,<)\) is potentially isomorphic to itself, \((\alpha,<)\models
\sigma\) and \(\sigma\in T(\alpha)\).  Since \((\beta,<)\) is not
potentially isomorphic to \((\alpha,<)\), \((\beta,<)\not\models\sigma\)
and \(\sigma\notin T(\beta)\).  So \(T(\alpha)\neq T(\beta)\).

\emph{\(A\) is a set.}  The class
\(B=\{X\in\mathcal P(S):X=T(\alpha)\text{ for some }\alpha\in A\}\) is a
set by Power Set and Separation.  By injectivity, each \(X\in B\) equals
\(T(\alpha)\) for exactly one \(\alpha\in A\); call it \(G(X)\).  The
class function \(G\) maps \(B\) onto \(A\), so \(A=G[B]\) is a set by
Replacement.

Let \(\xi\) be the supremum of the set \(\{1\}\cup\{\alpha+1:\alpha\in A\}\)
of ordinals.  If \(\alpha\geq\xi\), then \(\alpha\geq1\) and
\(\alpha\notin A\), since every member of \(A\) is below \(\xi\); so no
\(L_{\omega_1\omega}\)-sentence is equivalent to \(\delta_{(\alpha,<)}\).
Such an \(\alpha\) is uncountable: a countable \(\alpha\geq1\) lies in
\(A\) by clause~1, hence is below \(\xi\).

\emph{Different families of model classes.}  By Convention~\ref{def:p02-w02-p02-039-notation},
every \(L_{\omega_1\omega}\)-sentence is an \(L_{\infty\omega}\)-sentence
with the same models.  The \(L_{\infty\omega}\)-sentence
\(\delta_{(\xi,<)}\) has a model class that, by clause~2, belongs to no
\(L_{\omega_1\omega}\)-sentence.

\emph{The \(L_{\infty\omega}\)-sentences over \(\{<\}\) do not form a set.}
Suppose they formed a set \(S'\).  Let \(D\subseteq S'\) consist of those
\(\delta\) for which some ordinal \(\alpha\geq1\) has the property that the
models of \(\delta\) are exactly the structures potentially isomorphic to
\((\alpha,<)\); \(D\) is a set by Separation.  For \(\delta\in D\) that
ordinal is unique: if \(\alpha\) and \(\beta\) both have the property, then
\((\alpha,<)\models\delta\) because \((\alpha,<)\) is potentially
isomorphic to itself, so \((\alpha,<)\) is potentially isomorphic to
\((\beta,<)\), and \(\alpha=\beta\) by
Lemma~\ref{thm:p02-w02-p02-039-ordinal-rigidity}.  Call it \(H(\delta)\).
By Theorem~\ref{thm:p08-bvb501-potential-isomorphism}, every ordinal
\(\alpha\geq1\) equals \(H(\delta_{(\alpha,<)})\), so \(H\) maps the set
\(D\) onto the class of ordinals \(\alpha\geq1\).  Replacement would make
that class a set, and adding \(0\) would make the class of all ordinals a
set \(O\).  That is impossible: \(O\) would be a transitive set well-ordered
by membership, hence an ordinal, so \(O\in O\), while no ordinal is a member
of itself.  This contradiction completes the proof.
\end{proof}
\end{annexitem}
